\documentclass[aspectratio=1610]{beamer}
\usepackage[T1]{fontenc}
\usetheme{wildcat}
\usetikzlibrary{arrows.meta,angles,quotes,calc,intersections,positioning}

\usepackage{amsmath,amssymb,amsfonts}
\usepackage{booktabs}
\usepackage{relsize}
\usepackage{pgfplots}
\pgfplotsset{compat=1.16}
\usepackage{array}
\usepackage{siunitx}
\usepackage{cancel}
\usepackage{jupyter}
\usepackage{minted}
\usepackage{xfrac}
\usepackage{mathtools}

\let\oldfootnotesize\footnotesize
\renewcommand*{\footnotesize}{\oldfootnotesize\tiny}

\def\mathdefault#1{#1}
\everymath=\expandafter{\the\everymath\displaystyle}


\title{Elements of Modern C++ and Fortran \\
       {\small\it ME 701 - Phase III - Lesson 01}}

\date{\input{term.txt} \\ {\footnotesize Git SHA: \input{git_sha.txt}}}

\author{Jeremy Roberts}


\definecolor{ksupurple}{HTML}{512888}
\definecolor{orange}{HTML}{CA7C1B}
\definecolor{skyblue}{RGB}{180,220,255}
\definecolor{light-gray}{gray}{0.97}

\usepackage{listings}

\lstset{basicstyle=\fontsize{7}{8}\selectfont\ttfamily,
        numbers=left,
        numberstyle=\fontsize{5}{6}\selectfont\ttfamily,
        numbersep=5pt,                  
        backgroundcolor=\color{light-gray},
        frame=single, 
        %rulesepcolor=\color{red},
        keywordstyle=\color[rgb]{0,0,1},
        commentstyle=\color[rgb]{0.133,0.545,0.133},
        stringstyle=\color[rgb]{0.627,0.126,0.941},
        captionpos=b,
        title=\lstname,
        showstringspaces=false
       }


\begin{document}

\begin{frame}
\titlepage
\end{frame}
 
 
%%%%%%%%%%%%%%%%%%%%%%%%%%%%%%%%%%%%%%%%%%%%%%%%%%%
\begin{frame}{Primary Objective}

Students will be able to 

\vfill

\begin{quote}
\textcolor{wcprimary}{recognize, understand, and apply basic programming patterns for technical computing using the C++ and Modern Fortran languages.
}
\end{quote}

\vfill 

\end{frame}


%%%%%%%%%%%%%%%%%%%%%%%%%%%%%%%%%%%%%%%%%%%%%%%%%%%
\begin{frame}{References}

\begin{enumerate}
 \item Metcalf, Michael, et al. {\bf Modern Fortran Explained: Incorporating Fortran 2023}. Oxford University Press, 2024.
 \item Stroustrup, Bjarne. {\bf A Tour of C++}. Addison-Wesley Professional, 2022.
 \item {\bf AMTH 483/583 High-Performance Scientific Computing} (University of Washington), \url{https://amath583.github.io/sp22/index.html}
 \item {\bf Modern Fortran} (West Virginia University), \url{https://wvuhpc.github.io/Modern-Fortran/}
\end{enumerate}

\end{frame}


\section{C++ and Modern Fortran Overview}

%------------------------------------------------------------------------------%
\begin{frame}{Hello World}

These and other examples are in \textcolor{wcprimary}{\tt me701/examples/compiled}.


\begin{columns}[c]
  \begin{column}{0.5\textwidth}
    \lstinputlisting[language=C++]{compiled_code/hello_world.cc}
  \end{column}
  \begin{column}{0.5\textwidth}
    \lstinputlisting[language=Fortran]{compiled_code/hello_world.f90}
  \end{column}
\end{columns}

\end{frame}

%------------------------------------------------------------------------------%
\begin{frame}{VS Code Extensions for C++ and Fortran}

For C++:
\begin{itemize}
 \item C/C++ (Microsoft)
 \item C/C++ Extension Pack (Microsoft)
 \item C/C++ Themes (Microsoft)
\end{itemize}

\vfill 

For Fortran:
\begin{itemize}
 \item Modern Fortran (fortran-lang.org; unsigned, so need to click gear to install)
 \item {\tt pip install fortls} (in your ME 701 environment)
\end{itemize}

\end{frame}

%------------------------------------------------------------------------------%
\begin{frame}{Compiling Your First Program}

For C++, use (in the command line)
\begin{equation*}
 \overbrace{\tt g\!+\!+}^{\text{compiler}}~
   \underbrace{\tt{hello\_world.cc}}_{\text{file to compile}}~
     \overbrace{\tt{-o}}^{\text{output as}}~
       \underbrace{\tt{hello\_world}}_{\text{this executable}} 
\end{equation*}
\vfill 

For Fortran, use
\begin{equation*}
 \overbrace{\tt gfortran}^{\text{compiler}}~
   \underbrace{\tt{hello\_world.f90}}_{\text{file to compile}}~
     \overbrace{\tt{-o}}^{\text{output as}}~
       \underbrace{\tt{hello\_world}}_{\text{this executable}} 
\end{equation*}

\vfill 
Use \textcolor{wcprimary}{{\tt sudo apt-get install g++ gfortran}} to get them.
{\bf Now try them!}

\end{frame}

%------------------------------------------------------------------------------%
\begin{frame}{Compiler Options}

{\tt g++} and {\tt gfortran} are part of the GNU compiler set and share 
several key compiler options that may (or may not) work with compilers from
other vendors; these include:
\begin{itemize}
 \item {\tt -Wall} -- warn us of anything unexpected but make the executable
 \item {\tt -Werror} -- turn any warning into an error
 \item {\tt -O} -- (that's an ``Oh'') use optimization (or {\tt -ON} for $N=0,1,2,3$ for 
       various levels of optimization)
 \item {\tt -g} -- produce debugging information
 \item {\tt -pg} -- produce profiling information
\end{itemize}


\end{frame}

\begin{frame}{Declaring Variables}
\begin{columns}[c]
  \begin{column}{0.5\textwidth}
    \lstinputlisting[language=C++]{compiled_code/declaring.cc}
  \end{column}
  \begin{column}{0.5\textwidth}
    \lstinputlisting[language=Fortran]{compiled_code/declaring.f90}
  \end{column}
\end{columns}
\end{frame}

\begin{frame}{Simple Math}
\begin{columns}[c]
  \begin{column}{0.5\textwidth}
    \lstinputlisting[language=C++]{compiled_code/simple_math.cc}
  \end{column}
  \begin{column}{0.5\textwidth}
    \lstinputlisting[language=Fortran]{compiled_code/simple_math.f90}
  \end{column}
\end{columns}
\end{frame}

\begin{frame}{Control of Program Flow -- If's}
\begin{columns}[c]
  \begin{column}{0.5\textwidth}
    \lstinputlisting[language=C++]{compiled_code/control.cc}
  \end{column}
  \begin{column}{0.5\textwidth}
    \lstinputlisting[language=Fortran]{compiled_code/control.f90}
  \end{column}
\end{columns}
\end{frame}

\begin{frame}{Control of Program Flow -- Switches}
\begin{columns}[c]
  \begin{column}{0.5\textwidth}
    \lstinputlisting[language=C++]{compiled_code/control2.cc}
  \end{column}
  \begin{column}{0.5\textwidth}
    \lstinputlisting[language=Fortran]{compiled_code/control2.f90}
  \end{column}
\end{columns}
\end{frame}

\begin{frame}{Loops}
\begin{columns}[c]
  \begin{column}{0.5\textwidth}
    \lstinputlisting[language=C++]{compiled_code/loops.cc}
  \end{column}
  \begin{column}{0.5\textwidth}
    \lstinputlisting[language=Fortran]{compiled_code/loops.f90}
  \end{column}
\end{columns}
\end{frame}

\begin{frame}{Functions}
\begin{columns}[c]
  \begin{column}{0.5\textwidth}
    \lstinputlisting[language=C++]{compiled_code/functions.cc}
  \end{column}
  \begin{column}{0.5\textwidth}
    \lstinputlisting[language=Fortran]{compiled_code/functions.f90}
  \end{column}
\end{columns}
\end{frame}




%------------------------------------------------------------------------------%
\begin{frame}{Command Line Arguments}
\begin{columns}[c]
  \begin{column}{0.5\textwidth}
    \lstinputlisting[language=C++]{compiled_code/cl.cc}
  \end{column}
  \begin{column}{0.5\textwidth}
    \lstinputlisting[language=Fortran]{compiled_code/cl.f90}
  \end{column}
\end{columns}
\end{frame}


%------------------------------------------------------------------------------%
\begin{frame}{File I/O}
\begin{columns}[c]
  \begin{column}{0.5\textwidth}
    \lstinputlisting[language=C++]{compiled_code/file_io.cc}
  \end{column}
  \begin{column}{0.5\textwidth}
    \lstinputlisting[language=Fortran]{compiled_code/file_io.f90}
  \end{column}
\end{columns}
\end{frame}

%------------------------------------------------------------------------------%
\begin{frame}{File I/O}
\begin{columns}[c]
  \begin{column}{0.5\textwidth}
    \lstinputlisting[language=C++]{compiled_code/file_io.cc}
  \end{column}
  \begin{column}{0.5\textwidth}
    \lstinputlisting[language=Fortran]{compiled_code/file_io.f90}
  \end{column}
\end{columns}
\end{frame}

%==============================================================================%
\section{C++ Arrays}
%==============================================================================%

%------------------------------------------------------------------------------%
\begin{frame}[allowframebreaks]{C++ Arrays - Main Program}
 \lstinputlisting[language=C++]{compiled_code/basic_arrays.cc}
\end{frame}

%------------------------------------------------------------------------------%
\begin{frame}[allowframebreaks]{C++ Arrays - Function Header}
 \lstinputlisting[language=C++]{compiled_code/basic_arrays_functions.hh}
\end{frame}

%------------------------------------------------------------------------------%
\begin{frame}[allowframebreaks]{C++ Arrays - Function Definitions}
 \lstinputlisting[language=C++]{compiled_code/basic_arrays_functions.cc}
\end{frame}

%------------------------------------------------------------------------------%
\begin{frame}{Compiling}

 Go ahead, try {\tt g++ basic\_array.cc}.

 \pause
 
 \vfill
 The error is a {\bf linking error}.  
 
 \pause 
 
 \vfill
 Comment out the {\tt \#include "basic\_arrays\_functions.hh"} line 
 in {\tt basic\_array.cc} and try compiling again.
 
 \pause 
 \vfill
 The error is a {\bf syntactical error}---all names, including 
 functions, need to be defined before use.
 
 \pause 
 \vfill
 The right way (after uncommenting the {\tt .hh} file): \\
 
 {\tt g++ basic\_array\_functions.cc basic\_array.cc}\\
 
 (where the order matters!)
 
 
\end{frame}

%------------------------------------------------------------------------------%
\begin{frame}{Compiling with {\tt make}}
  \lstinputlisting[language=make]{compiled_code/make_basic_arrays}
\end{frame}

%------------------------------------------------------------------------------%
\begin{frame}[allowframebreaks]{Fortran Arrays - Main Program}
 \lstinputlisting[language=Fortran]{compiled_code/basic_arrays.f90}
\end{frame}

%------------------------------------------------------------------------------%
\begin{frame}[allowframebreaks]{Fortran Arrays - Functions Module}
 \lstinputlisting[language=Fortran]{compiled_code/basic_arrays_module.f90}
\end{frame}



\end{document}

